% Setting up the document class and essential packages
\documentclass[a4paper,12pt]{article}
\usepackage{amsmath, amsthm, amssymb,color,braket}
\usepackage{geometry}
\geometry{margin=0.5in} % Set 1.5-inch margins on all sides

\usepackage{parskip}
\renewcommand{\familydefault}{\sfdefault}


\newtheoremstyle{mystyle}%                % Name
{}%                                     % Space above
{}%                                     % Space below
{\itshape}%                             % Body font
{}%                                     % Indent amount
{\bfseries}%                            % Theorem head font
{.}%                                    % Punctuation after theorem head
{ }%                                    % Space after theorem head, ' ', or \newline
{}%     


\theoremstyle{mystyle}
% Configuring theorem-like environments
\newtheorem{theorem}{Theorem}[section]
\newtheorem{lemma}[theorem]{Lemma}
\newtheorem{corollary}[theorem]{Corollary}

\theoremstyle{mystyle}
\newtheorem{proposition}{Proposition}
\newtheorem{definition}{\color{blue} Definition}
\newtheorem{remark}{\color{red} Remark}
\newtheorem{example}{Example}

% Setting up section numbering
\numberwithin{equation}{section}
\title{Algebraic Definition of Tangent Structure on Manifolds}
\begin{document}
	\maketitle
	Algebraic definition of differential forms relies on the definition of sheaf.The concept of Sheaf is used to describe how the information of a local structure is integrated to a global structure.
	\begin{definition}
		A sheaf on a topological space $(M,\tau)$ is a map $F$ that assign each open set $U$ a abelian group $F(U)$ and for any $U\subseteq V$,the inclusion must give a homeomorphism of the abelian group
		$$\rho_{UV}:F(V)\to F(U)$$
		and that
		\\
		(1):$\rho$ must be transitive (i.e. $[0]$ for 1-cycles):$\rho_{UV}\circ \rho_{VW}=\rho_{UW}$.
		\\
		It is called presheaf if (1) is satisfied.
		\\
		(2):The following short sequence is exact
		$$0\to F(M)\to \prod_{\alpha} F(U_{\alpha})\to \prod_{\alpha,\beta} F(U_{\alpha}\cap U_\beta)\to 0$$
		The first map is simply the decomposition map $\oplus_{\alpha} \rho_{U_\alpha M}$,and the second map is the 2-cocycle map
		$$\oplus_{\alpha,\beta}(\rho_{U_\alpha\cap U_\beta,U_\alpha}-\rho_{U_\alpha\cap U_\beta,U_\beta})$$
	\end{definition}
	\begin{example}
		$O_M$,Smooth function sheaf on $C^\infty$ manifold $M$.$F(U)$ is simply all the smooth functions on $U$ equipped with addition.The restriction map is simply the restriction map of functions.
	\end{example}
	A sheaf is a ring sheaf if every $F(U)$ is a ring.
	\\ 
	Using $C^\infty$ function sheaf we can define the differential sheaf.Consider the diagonal map on manifold $M$
	$$\Delta:M\to M\times M,p\to (p,p)$$
	Take any open set $U$ in which we have a local coordinate $x^i$.Then we have local coordinates in $U\times U$
	$$x^i\otimes 1,1\otimes x^i$$
	\begin{proposition}
		$\Delta$ induces a ring sheaf hom in the following sense:take any $f(p,q)\in O_{M\times M}$,define 
		$$\Delta^*(f)=f(p,p)\in O_{\Delta(M)}$$
	\end{proposition}
	\begin{proof}
		$\Delta^*$ is apparently a ring hom.The only thing we need to verify is that restriction commute with hom.This is also trivial because
		$$f|_{U\times U}(p,p)=\Delta^*f|_U(p)$$
	\end{proof}
	\begin{proposition}
	Take any open $U\times U$ with the above local coordinates.Then the ideal $I=ker(\Delta^*)$ is finitely generated by $x^i\otimes 1-1\otimes x^i$.
	\end{proposition}
	\begin{proof}
		It suffices to prove $\forall f(p,p)=0,f\in C^\infty,f(p,q)=(\sum_i^n \lambda_i(p^i\otimes 1-1\otimes q^i)g_i)$.WLOG,assume $U$ is convex.By fundamental theorem of calculus,
		$$f(p,q)=f(p,p)+\int_0^1 (p\otimes 1-1\otimes q)\cdot \nabla f(p,p+s(q-p))ds$$
		Since $f(p,p)=0$,we see that indeed the ideal is generated by $p^i\otimes 1-1\otimes q^i$.
	\end{proof}
	\begin{definition}
		Let $I$ be the ideal defined in proposition 2.Consider the quotient ring $I/I^2$.$I/I^2$ naturally has a $O_{M\times M}$-module structure given by ring multiplication.Consider the pull back module $\Delta^*(I/I^2)$ (notice that we are not considering pull back ring!!!If so $\Delta^*(I/I^2)$ would be zero) such that
		$$f\cdot \Delta^*(I/I^2)=(\Delta_*f)\cdot I/I^2,f\in O_{M}$$
		The pull back is well defined (remember that push forward of functions is not uniquely defined.) because take any $\bar{f},\bar{g}\in \Delta_*f$,we must have $\bar{f}-\bar{g}\in I$.Then
		$$(\bar{f}-\bar{g})\cdot I/I^2\in I\cdot I/I^2=0$$
		So the definition is independent of representative.The pull back $O_M$-module $\Delta^*(I/I^2)$ is defined as $\Omega_M^1$,the differential sheaf on $M$,with restriction maps being restriction of functions.
	\end{definition}
	\begin{remark}
		Since we have already proved locally the ideal $I$ is generated by $x^i\otimes 1-1\otimes x^i$,then the differential sheaf is locally generated by
		$$dx^i=\overline{x^i\otimes 1-1\otimes x^i}$$
		with relations $dx^idx^j=0,\forall i,j$.
	\end{remark}
	\begin{remark}
		There is a sheaf hom $d:O_M\to \Omega_M^1$ defined as follows.Define firstly a ring hom
		$$\delta:O_M\to O_{M\times M},f(x)\to f(p,q)=f(p)-f(q)$$
		And apparently $im(\delta)\in I$.So let
		$$d(f)=\Delta^*\pi \delta(f)$$
		Where $\pi:I\to I/I^2$ is the quotient map.This is clearly a sheaf hom(it is easy to verify all of this when the restriction map on sheaf is simply function restriction).We can now check the Leibniz rule.
		\\
		$\forall f,g\in O_M$,
		$$\delta(fg)(p,q)=f(p)g(p)-f(q)g(q)$$
		Since $\delta(f),\delta(g)\in I\implies \delta(f)\delta(g)\in I^2$,It is easy to notice that
		$$\pi\delta(fg)=\overline{\delta(fg)}=\overline{f(p)g(p)-f(q)g(q)-\delta(f)\delta(g)}=\overline{f(q)\delta(g)+g(q)\delta(f)}$$
		So
		$$\Delta^*\pi \delta(fg)=f\Delta^*\pi \delta(g)+g\Delta^*\pi \delta(f)=fd(g)+gd(f)$$
	\end{remark}
	Since we already know the differential sheaf is locally generated by $dx^i$,the tangent sheaf should be its dual (as we know in differential geometry).Then
	\begin{definition}
		The tangent sheaf on $M$ is
		$$\mathcal{T}_M=Hom(\Omega_M^1,O_M)$$
		with restriction
		$$\mathcal{T}_M|_{U}=Hom(\Omega_M^1|_{U},O_M|_{U})$$
	\end{definition}
\end{document}